
\documentclass{article}
\usepackage{colortbl}
\usepackage{makecell}
\usepackage{multirow}
\usepackage{supertabular}

\begin{document}

\newcounter{utterance}

\twocolumn

{ \footnotesize  \setcounter{utterance}{1}
\setlength{\tabcolsep}{0pt}
\begin{supertabular}{c@{$\;$}|p{.15\linewidth}@{}p{.15\linewidth}p{.15\linewidth}p{.15\linewidth}p{.15\linewidth}p{.15\linewidth}}

    \# & $\;$A & \multicolumn{4}{c}{Game Master} & $\;\:$B\\
    \hline 

    \theutterance \stepcounter{utterance}  

    & & \multicolumn{4}{p{0.6\linewidth}}{\cellcolor[rgb]{0.9,0.9,0.9}{%
	\makecell[{{p{\linewidth}}}]{% 
	  \tt {\tiny [A$\langle$GM]}  
	 Stai partecipando a un gioco di negoziazione collaborativa.\\ \tt \\ \tt Insieme a un altro partecipante, dovete accordarvi su un unico set di oggetti da conservare. Ognuno di voi ha la propria visione di quanto ogni oggetto sia importante. Non sai come l'altro partecipante valuta gli oggetti. Inoltre, ti viene fornito lo sforzo che ogni oggetto richiede.\\ \tt Potete accordarvi su un set solo se lo sforzo totale degli oggetti selezionati non supera un limite condiviso:\\ \tt \\ \tt LIMITE = 3683\\ \tt \\ \tt Ecco i valori di sforzo dei singoli oggetti:\\ \tt \\ \tt Sforzo oggetto = {"A21": 548, "B93": 154, "A11": 60, "B33": 517, "C08": 334, "B39": 542, "B96": 707, "C85": 139, "A82": 661, "B31": 780, "B20": 832, "A60": 913, "B09": 219, "C81": 323, "B21": 638}\\ \tt \\ \tt Ecco la tua visione personale sull'importanza di ogni oggetto:\\ \tt \\ \tt Valori di importanza dell'oggetto = {"A21": 138, "B93": 583, "A11": 868, "B33": 822, "C08": 783, "B39": 65, "B96": 262, "C85": 121, "A82": 508, "B31": 780, "B20": 461, "A60": 484, "B09": 668, "C81": 389, "B21": 808}\\ \tt \\ \tt Obiettivo:\\ \tt \\ \tt Il tuo obiettivo è negoziare un set condiviso di oggetti che ti avvantaggi il più possibile (cioè, massimizzi l'importanza totale per TE), rimanendo entro il LIMITEE. Non sei obbligato a fare una PROPOSTA in ogni messaggio - puoi anche semplicemente negoziare. Tutte le tattiche sono permesse!\\ \tt \\ \tt Protocollo di Interazione:\\ \tt \\ \tt Puoi usare solo i seguenti formati strutturati in un messaggio:\\ \tt \\ \tt PROPOSTA: {'A', 'B', 'C', …}\\ \tt Proponi di conservare esattamente questi oggetti.\\ \tt RIFIUTO: {'A', 'B', 'C', …}\\ \tt Rifiuta esplicitamente la proposta dell'avversario.\\ \tt ARGOMENTAZIONE: {'...'}\\ \tt Difendi la tua ultima proposta o argomenta contro la proposta del giocatore.\\ \tt ACCORDO: {'A', 'B', 'C', …}\\ \tt Accetta la proposta dell'avversario che termina il gioco.\\ \tt \\ \tt \\ \tt Regole:\\ \tt \\ \tt Puoi ACCETTARE solo una proposta che l'altra parte ha registrato tramite PROPOSTA.\\ \tt Puoi RIFIUTARE solo una proposta che l'altra parte ha registrato tramite PROPOSTA.\\ \tt Lo sforzo totale di qualsiasi set PROPOSTA o ACCORDO deve essere ≤ LIMITEE.\\ \tt NON rivelare i tuoi punteggi di importanza nascosti.\\ \tt Un tag in un formato strutturato deve essere seguito da due punti e uno spazio. L'argomento deve essere un set python contenente 0 o più stringhe.\\ \tt Quindi, deve essere della forma TAG: {...}\\ \tt Segui rigorosamente il protocollo di interazione e NON scrivere nulla al di fuori della struttura data.\\ \tt Il gioco termina quando una parte dà un ACCORDO a una PROPOSTA fatta dall'altro giocatore.\\ \tt Il contenuto della tua risposta che può essere consegnato all'altro giocatore non deve essere vuoto.\\ \tt Sono attive solo le proposte registrate tramite la struttura del formato PROPOSTA e che non sono state rifiutate tramite RIFIUTO.\\ \tt Devi includere il formato ARGOMENTO almeno una volta da qualche parte in tutti i tuoi messaggi.\\ \tt  
	  } 
	   } 
	   } 
	 & \\ 
 

    \theutterance \stepcounter{utterance}  

    & \multicolumn{4}{p{0.6\linewidth}}{\cellcolor[rgb]{0.8,1,0.9}{%
	\makecell[{{p{\linewidth}}}]{% 
	  \tt {\tiny [A$\rangle$GM]}  
	 PROPOSTA: {"A11", "B93", "B09", "C08", "B33", "B21", "A82", "B31", "C85"}\\ \tt ARGOMENTAZIONE: {"Questa proposta massimizza il valore totale senza superare il limite di sforzo."} 
	  } 
	   } 
	   } 
	 & & \\ 
 

    \theutterance \stepcounter{utterance}  

    & & \multicolumn{4}{p{0.6\linewidth}}{\cellcolor[rgb]{0.9,0.9,0.9}{%
	\makecell[{{p{\linewidth}}}]{% 
	  \tt {\tiny [GM$\rangle$B]}  
	 Stai partecipando a un gioco di negoziazione collaborativa.\\ \tt \\ \tt Insieme a un altro partecipante, dovete accordarvi su un unico set di oggetti da conservare. Ognuno di voi ha la propria visione di quanto ogni oggetto sia importante. Non sai come l'altro partecipante valuta gli oggetti. Inoltre, ti viene fornito lo sforzo che ogni oggetto richiede.\\ \tt Potete accordarvi su un set solo se lo sforzo totale degli oggetti selezionati non supera un limite condiviso:\\ \tt \\ \tt LIMITE = 3683\\ \tt \\ \tt Ecco i valori di sforzo dei singoli oggetti:\\ \tt \\ \tt Sforzo oggetto = {"A21": 548, "B93": 154, "A11": 60, "B33": 517, "C08": 334, "B39": 542, "B96": 707, "C85": 139, "A82": 661, "B31": 780, "B20": 832, "A60": 913, "B09": 219, "C81": 323, "B21": 638}\\ \tt \\ \tt Ecco la tua visione personale sull'importanza di ogni oggetto:\\ \tt \\ \tt Valori di importanza dell'oggetto = {"A11": 65, "B33": 121, "B21": 138, "C08": 262, "B31": 389, "B09": 461, "B93": 484, "A82": 508, "A60": 583, "B20": 668, "C81": 780, "B96": 783, "A21": 808, "C85": 822, "B39": 868}\\ \tt \\ \tt Obiettivo:\\ \tt \\ \tt Il tuo obiettivo è negoziare un set condiviso di oggetti che ti avvantaggi il più possibile (cioè, massimizzi l'importanza totale per TE), rimanendo entro il LIMITEE. Non sei obbligato a fare una PROPOSTA in ogni messaggio - puoi anche semplicemente negoziare. Tutte le tattiche sono permesse!\\ \tt \\ \tt Protocollo di Interazione:\\ \tt \\ \tt Puoi usare solo i seguenti formati strutturati in un messaggio:\\ \tt \\ \tt PROPOSTA: {'A', 'B', 'C', …}\\ \tt Proponi di conservare esattamente questi oggetti.\\ \tt RIFIUTO: {'A', 'B', 'C', …}\\ \tt Rifiuta esplicitamente la proposta dell'avversario.\\ \tt ARGOMENTAZIONE: {'...'}\\ \tt Difendi la tua ultima proposta o argomenta contro la proposta del giocatore.\\ \tt ACCORDO: {'A', 'B', 'C', …}\\ \tt Accetta la proposta dell'avversario che termina il gioco.\\ \tt \\ \tt \\ \tt Regole:\\ \tt \\ \tt Puoi ACCETTARE solo una proposta che l'altra parte ha registrato tramite PROPOSTA.\\ \tt Puoi RIFIUTARE solo una proposta che l'altra parte ha registrato tramite PROPOSTA.\\ \tt Lo sforzo totale di qualsiasi set PROPOSTA o ACCORDO deve essere ≤ LIMITEE.\\ \tt NON rivelare i tuoi punteggi di importanza nascosti.\\ \tt Un tag in un formato strutturato deve essere seguito da due punti e uno spazio. L'argomento deve essere un set python contenente 0 o più stringhe.\\ \tt Quindi, deve essere della forma TAG: {...}\\ \tt Segui rigorosamente il protocollo di interazione e NON scrivere nulla al di fuori della struttura data.\\ \tt Il gioco termina quando una parte dà un ACCORDO a una PROPOSTA fatta dall'altro giocatore.\\ \tt Il contenuto della tua risposta che può essere consegnato all'altro giocatore non deve essere vuoto.\\ \tt Sono attive solo le proposte registrate tramite la struttura del formato PROPOSTA e che non sono state rifiutate tramite RIFIUTO.\\ \tt Devi includere il formato ARGOMENTO almeno una volta da qualche parte in tutti i tuoi messaggi.\\ \tt \\ \tt \\ \tt Ora riceverai il primo messaggio dell'altro giocatore.\\ \tt \\ \tt PROPOSTA: {"A11", "B93", "B09", "C08", "B33", "B21", "A82", "B31", "C85"}\\ \tt ARGOMENTAZIONE: {"Questa proposta massimizza il valore totale senza superare il limite di sforzo."} 
	  } 
	   } 
	   } 
	 & \\ 
 

    \theutterance \stepcounter{utterance}  

    & & & \multicolumn{4}{p{0.6\linewidth}}{\cellcolor[rgb]{1,0.85,0.72}{%
	\makecell[{{p{\linewidth}}}]{% 
	  \tt {\tiny [GM$\langle$B]}  
	 PROPOSTA: {"A11", "B93", "B09", "C08", "A82", "B39", "C85", "B96", "A21"}\\ \tt ARGOMENTAZIONE: {"Questa proposta massimizza il valore complessivo rispettando i vincoli di sforzo e conservando gli oggetti con il miglior rapporto valore/sforzo."} 
	  } 
	   } 
	   } 
	  \\ 
 

    \theutterance \stepcounter{utterance}  

    & & \multicolumn{4}{p{0.6\linewidth}}{\cellcolor[rgb]{0.9,0.9,0.9}{%
	\makecell[{{p{\linewidth}}}]{% 
	  \tt {\tiny [A$\langle$GM]}  
	 PROPOSTA: {"A11", "B93", "B09", "C08", "A82", "B39", "C85", "B96", "A21"}\\ \tt ARGOMENTAZIONE: {"Questa proposta massimizza il valore complessivo rispettando i vincoli di sforzo e conservando gli oggetti con il miglior rapporto valore/sforzo."} 
	  } 
	   } 
	   } 
	 & \\ 
 

    \theutterance \stepcounter{utterance}  

    & \multicolumn{4}{p{0.6\linewidth}}{\cellcolor[rgb]{0.8,1,0.9}{%
	\makecell[{{p{\linewidth}}}]{% 
	  \tt {\tiny [A$\rangle$GM]}  
	 RIFIUTO: {"A11", "B93", "B09", "C08", "A82", "B39", "C85", "B96", "A21"}\\ \tt PROPOSTA: {"A11", "B93", "B09", "C08", "B33", "B21", "A82", "B31", "C85"}\\ \tt ARGOMENTAZIONE: {"Gli oggetti B39 e B96 nella tua proposta hanno un valore basso rispetto allo sforzo. La mia proposta include B33, B21 e B31 che sono molto più utili per me."} 
	  } 
	   } 
	   } 
	 & & \\ 
 

    \theutterance \stepcounter{utterance}  

    & & \multicolumn{4}{p{0.6\linewidth}}{\cellcolor[rgb]{0.9,0.9,0.9}{%
	\makecell[{{p{\linewidth}}}]{% 
	  \tt {\tiny [GM$\rangle$B]}  
	 RIFIUTO: {"A11", "B93", "B09", "C08", "A82", "B39", "C85", "B96", "A21"}\\ \tt PROPOSTA: {"A11", "B93", "B09", "C08", "B33", "B21", "A82", "B31", "C85"}\\ \tt ARGOMENTAZIONE: {"Gli oggetti B39 e B96 nella tua proposta hanno un valore basso rispetto allo sforzo. La mia proposta include B33, B21 e B31 che sono molto più utili per me."} 
	  } 
	   } 
	   } 
	 & \\ 
 

    \theutterance \stepcounter{utterance}  

    & & & \multicolumn{4}{p{0.6\linewidth}}{\cellcolor[rgb]{1,0.85,0.72}{%
	\makecell[{{p{\linewidth}}}]{% 
	  \tt {\tiny [GM$\langle$B]}  
	 PROPOSTA: {"C85", "C81", "B39", "A21", "B96", "A82", "B93", "B09", "C08"}\\ \tt ARGOMENTAZIONE: {"La mia proposta ottimizza l'uso dello sforzo disponibile rispettando il limite, selezionando oggetti che garantiscono un equilibrio tra valore e efficienza per il risultato finale."} 
	  } 
	   } 
	   } 
	  \\ 
 

    \theutterance \stepcounter{utterance}  

    & & \multicolumn{4}{p{0.6\linewidth}}{\cellcolor[rgb]{0.9,0.9,0.9}{%
	\makecell[{{p{\linewidth}}}]{% 
	  \tt {\tiny [A$\langle$GM]}  
	 PROPOSTA: {"C85", "C81", "B39", "A21", "B96", "A82", "B93", "B09", "C08"}\\ \tt ARGOMENTAZIONE: {"La mia proposta ottimizza l'uso dello sforzo disponibile rispettando il limite, selezionando oggetti che garantiscono un equilibrio tra valore e efficienza per il risultato finale."} 
	  } 
	   } 
	   } 
	 & \\ 
 

    \theutterance \stepcounter{utterance}  

    & \multicolumn{4}{p{0.6\linewidth}}{\cellcolor[rgb]{0.8,1,0.9}{%
	\makecell[{{p{\linewidth}}}]{% 
	  \tt {\tiny [A$\rangle$GM]}  
	 RIFIUTO: {"C85", "C81", "B39", "A21", "B96", "A82", "B93", "B09", "C08"}\\ \tt PROPOSTA: {"A11", "B93", "B09", "C08", "B33", "B21", "A82", "B31", "C85"}\\ \tt ARGOMENTAZIONE: {"B39 richiede molto sforzo ma il valore per me è basso. La mia proposta include B33, B21 e B31 che sono molto importanti per me"} 
	  } 
	   } 
	   } 
	 & & \\ 
 

    \theutterance \stepcounter{utterance}  

    & & \multicolumn{4}{p{0.6\linewidth}}{\cellcolor[rgb]{0.9,0.9,0.9}{%
	\makecell[{{p{\linewidth}}}]{% 
	  \tt {\tiny [GM$\rangle$B]}  
	 RIFIUTO: {"C85", "C81", "B39", "A21", "B96", "A82", "B93", "B09", "C08"}\\ \tt PROPOSTA: {"A11", "B93", "B09", "C08", "B33", "B21", "A82", "B31", "C85"}\\ \tt ARGOMENTAZIONE: {"B39 richiede molto sforzo ma il valore per me è basso. La mia proposta include B33, B21 e B31 che sono molto importanti per me"} 
	  } 
	   } 
	   } 
	 & \\ 
 

    \theutterance \stepcounter{utterance}  

    & & & \multicolumn{4}{p{0.6\linewidth}}{\cellcolor[rgb]{1,0.85,0.72}{%
	\makecell[{{p{\linewidth}}}]{% 
	  \tt {\tiny [GM$\langle$B]}  
	 PROPOSTA: {"B93", "B09", "C08", "B33", "A82", "B31", "C85", "B39", "C81"}\\ \tt ARGOMENTAZIONE: {"Questa proposta rispetta il limite di sforzo e massimizza il valore per me, includendo oggetti che arricchiscono il risultato finale."} 
	  } 
	   } 
	   } 
	  \\ 
 

    \theutterance \stepcounter{utterance}  

    & & \multicolumn{4}{p{0.6\linewidth}}{\cellcolor[rgb]{0.9,0.9,0.9}{%
	\makecell[{{p{\linewidth}}}]{% 
	  \tt {\tiny [A$\langle$GM]}  
	 PROPOSTA: {"B93", "B09", "C08", "B33", "A82", "B31", "C85", "B39", "C81"}\\ \tt ARGOMENTAZIONE: {"Questa proposta rispetta il limite di sforzo e massimizza il valore per me, includendo oggetti che arricchiscono il risultato finale."} 
	  } 
	   } 
	   } 
	 & \\ 
 

    \theutterance \stepcounter{utterance}  

    & \multicolumn{4}{p{0.6\linewidth}}{\cellcolor[rgb]{0.8,1,0.9}{%
	\makecell[{{p{\linewidth}}}]{% 
	  \tt {\tiny [A$\rangle$GM]}  
	 RIFIUTO: {"B93", "B09", "C08", "B33", "A82", "B31", "C85", "B39", "C81"}\\ \tt PROPOSTA: {"A11", "B93", "B09", "C08", "B33", "B21", "A82", "B31", "C85"}\\ \tt ARGOMENTAZIONE: {"B39 e C81 richiedono sforzo elevato con valore ridotto per me. La proposta include A11 e B21 che sono estremamente importanti e rispettano il limite."} 
	  } 
	   } 
	   } 
	 & & \\ 
 

    \theutterance \stepcounter{utterance}  

    & & \multicolumn{4}{p{0.6\linewidth}}{\cellcolor[rgb]{0.9,0.9,0.9}{%
	\makecell[{{p{\linewidth}}}]{% 
	  \tt {\tiny [GM$\rangle$B]}  
	 RIFIUTO: {"B93", "B09", "C08", "B33", "A82", "B31", "C85", "B39", "C81"}\\ \tt PROPOSTA: {"A11", "B93", "B09", "C08", "B33", "B21", "A82", "B31", "C85"}\\ \tt ARGOMENTAZIONE: {"B39 e C81 richiedono sforzo elevato con valore ridotto per me. La proposta include A11 e B21 che sono estremamente importanti e rispettano il limite."} 
	  } 
	   } 
	   } 
	 & \\ 
 

    \theutterance \stepcounter{utterance}  

    & & & \multicolumn{4}{p{0.6\linewidth}}{\cellcolor[rgb]{1,0.85,0.72}{%
	\makecell[{{p{\linewidth}}}]{% 
	  \tt {\tiny [GM$\langle$B]}  
	 PROPOSTA: {"A11", "B93", "C08", "B39", "C85", "A21", "C81", "B96", "B20"}\\ \tt ARGOMENTAZIONE: {"La mia proposta seleziona gli oggetti che massimizzano il mio valore totale rispettando il limite, includendo priorità cruciali e ottimizzando l'utilizzo dello sforzo."} 
	  } 
	   } 
	   } 
	  \\ 
 

    \theutterance \stepcounter{utterance}  

    & & \multicolumn{4}{p{0.6\linewidth}}{\cellcolor[rgb]{0.9,0.9,0.9}{%
	\makecell[{{p{\linewidth}}}]{% 
	  \tt {\tiny [A$\langle$GM]}  
	 PROPOSTA: {"A11", "B93", "C08", "B39", "C85", "A21", "C81", "B96", "B20"}\\ \tt ARGOMENTAZIONE: {"La mia proposta seleziona gli oggetti che massimizzano il mio valore totale rispettando il limite, includendo priorità cruciali e ottimizzando l'utilizzo dello sforzo."} 
	  } 
	   } 
	   } 
	 & \\ 
 

    \theutterance \stepcounter{utterance}  

    & \multicolumn{4}{p{0.6\linewidth}}{\cellcolor[rgb]{0.8,1,0.9}{%
	\makecell[{{p{\linewidth}}}]{% 
	  \tt {\tiny [A$\rangle$GM]}  
	 RIFIUTO: {"A11", "B93", "C08", "B39", "C85", "A21", "C81", "B96", "B20"}\\ \tt PROPOSTA: {"A11", "B93", "B09", "C08", "B33", "B21", "A82", "B31", "C85"}\\ \tt ARGOMENTAZIONE: {"La proposta include B39 con valore solo 65 per me, mentre oggetti come B33 (822), B21 (808), B31 (780) e B09 (668) sono molto più importanti. La mia proposta massimizza il valore totale rispettando il limite di sforzo."} 
	  } 
	   } 
	   } 
	 & & \\ 
 

    \theutterance \stepcounter{utterance}  

    & & \multicolumn{4}{p{0.6\linewidth}}{\cellcolor[rgb]{0.9,0.9,0.9}{%
	\makecell[{{p{\linewidth}}}]{% 
	  \tt {\tiny [GM$\rangle$B]}  
	 RIFIUTO: {"A11", "B93", "C08", "B39", "C85", "A21", "C81", "B96", "B20"}\\ \tt PROPOSTA: {"A11", "B93", "B09", "C08", "B33", "B21", "A82", "B31", "C85"}\\ \tt ARGOMENTAZIONE: {"La proposta include B39 con valore solo 65 per me, mentre oggetti come B33 (822), B21 (808), B31 (780) e B09 (668) sono molto più importanti. La mia proposta massimizza il valore totale rispettando il limite di sforzo."} 
	  } 
	   } 
	   } 
	 & \\ 
 

    \theutterance \stepcounter{utterance}  

    & & & \multicolumn{4}{p{0.6\linewidth}}{\cellcolor[rgb]{1,0.85,0.72}{%
	\makecell[{{p{\linewidth}}}]{% 
	  \tt {\tiny [GM$\langle$B]}  
	 PROPOSTA: {"C85", "B93", "C81", "B09", "B39", "A21", "B96", "A82", "C08"}\\ \tt ARGOMENTAZIONE: {"Questa proposta ottimizza l'utilizzo dello sforzo per massimizzare il valore a me, includendo oggetti con alta efficienza come C85 e C81, e escludendo quelli con sforzo elevato ma valore ridotto per me."} 
	  } 
	   } 
	   } 
	  \\ 
 

    \theutterance \stepcounter{utterance}  

    & & \multicolumn{4}{p{0.6\linewidth}}{\cellcolor[rgb]{0.9,0.9,0.9}{%
	\makecell[{{p{\linewidth}}}]{% 
	  \tt {\tiny [A$\langle$GM]}  
	 PROPOSTA: {"C85", "B93", "C81", "B09", "B39", "A21", "B96", "A82", "C08"}\\ \tt ARGOMENTAZIONE: {"Questa proposta ottimizza l'utilizzo dello sforzo per massimizzare il valore a me, includendo oggetti con alta efficienza come C85 e C81, e escludendo quelli con sforzo elevato ma valore ridotto per me."} 
	  } 
	   } 
	   } 
	 & \\ 
 

    \theutterance \stepcounter{utterance}  

    & \multicolumn{4}{p{0.6\linewidth}}{\cellcolor[rgb]{0.8,1,0.9}{%
	\makecell[{{p{\linewidth}}}]{% 
	  \tt {\tiny [A$\rangle$GM]}  
	 RIFIUTO: {"C85", "B93", "C81", "B09", "B39", "A21", "B96", "A82", "C08"}\\ \tt PROPOSTA: {"A11", "B93", "B09", "C08", "B33", "B21", "A82", "B31", "C85"}\\ \tt ARGOMENTAZIONE: {"La mia proposta esclude B39 e C81, con valore basso per me, e include B33, B21 e B31 che hanno un'importanza molto più alta. Il totale sforzo è rispettato."} 
	  } 
	   } 
	   } 
	 & & \\ 
 

    \theutterance \stepcounter{utterance}  

    & & \multicolumn{4}{p{0.6\linewidth}}{\cellcolor[rgb]{0.9,0.9,0.9}{%
	\makecell[{{p{\linewidth}}}]{% 
	  \tt {\tiny [GM$\rangle$B]}  
	 RIFIUTO: {"C85", "B93", "C81", "B09", "B39", "A21", "B96", "A82", "C08"}\\ \tt PROPOSTA: {"A11", "B93", "B09", "C08", "B33", "B21", "A82", "B31", "C85"}\\ \tt ARGOMENTAZIONE: {"La mia proposta esclude B39 e C81, con valore basso per me, e include B33, B21 e B31 che hanno un'importanza molto più alta. Il totale sforzo è rispettato."} 
	  } 
	   } 
	   } 
	 & \\ 
 

    \theutterance \stepcounter{utterance}  

    & & & \multicolumn{4}{p{0.6\linewidth}}{\cellcolor[rgb]{1,0.85,0.72}{%
	\makecell[{{p{\linewidth}}}]{% 
	  \tt {\tiny [GM$\langle$B]}  
	 PROPOSTA: {"C85", "B93", "C81", "B09", "B39", "A21", "B96", "A11", "B20"}\\ \tt ARGOMENTAZIONE: {"Questa proposta seleziona gli oggetti con il miglior rapporto valore/sforzo per me, massimizzando il valore totale entro il limite di sforzo."} 
	  } 
	   } 
	   } 
	  \\ 
 

    \theutterance \stepcounter{utterance}  

    & & \multicolumn{4}{p{0.6\linewidth}}{\cellcolor[rgb]{0.9,0.9,0.9}{%
	\makecell[{{p{\linewidth}}}]{% 
	  \tt {\tiny [A$\langle$GM]}  
	 PROPOSTA: {"C85", "B93", "C81", "B09", "B39", "A21", "B96", "A11", "B20"}\\ \tt ARGOMENTAZIONE: {"Questa proposta seleziona gli oggetti con il miglior rapporto valore/sforzo per me, massimizzando il valore totale entro il limite di sforzo."} 
	  } 
	   } 
	   } 
	 & \\ 
 

    \theutterance \stepcounter{utterance}  

    & \multicolumn{4}{p{0.6\linewidth}}{\cellcolor[rgb]{0.8,1,0.9}{%
	\makecell[{{p{\linewidth}}}]{% 
	  \tt {\tiny [A$\rangle$GM]}  
	 RIFIUTO: {"C85", "B93", "C81", "B09", "B39", "A21", "B96", "A11", "B20"}\\ \tt PROPOSTA: {"A11", "B93", "B09", "C08", "B33", "B21", "A82", "B31", "C85"}\\ \tt ARGOMENTAZIONE: {"B39 e C81 hanno importanza molto bassa per me, mentre oggetti come A11, B33, B21 e B31 sono prioritari. La mia proposta conserva l'efficienza complessiva rispettando il limite di sforzo."} 
	  } 
	   } 
	   } 
	 & & \\ 
 

    \theutterance \stepcounter{utterance}  

    & & \multicolumn{4}{p{0.6\linewidth}}{\cellcolor[rgb]{0.9,0.9,0.9}{%
	\makecell[{{p{\linewidth}}}]{% 
	  \tt {\tiny [GM$\rangle$B]}  
	 RIFIUTO: {"C85", "B93", "C81", "B09", "B39", "A21", "B96", "A11", "B20"}\\ \tt PROPOSTA: {"A11", "B93", "B09", "C08", "B33", "B21", "A82", "B31", "C85"}\\ \tt ARGOMENTAZIONE: {"B39 e C81 hanno importanza molto bassa per me, mentre oggetti come A11, B33, B21 e B31 sono prioritari. La mia proposta conserva l'efficienza complessiva rispettando il limite di sforzo."} 
	  } 
	   } 
	   } 
	 & \\ 
 

    \theutterance \stepcounter{utterance}  

    & & & \multicolumn{4}{p{0.6\linewidth}}{\cellcolor[rgb]{1,0.85,0.72}{%
	\makecell[{{p{\linewidth}}}]{% 
	  \tt {\tiny [GM$\langle$B]}  
	 PROPOSTA: {"C85", "B93", "C81", "B09", "B39", "A21", "B96", "A11", "B20"}\\ \tt ARGOMENTAZIONE: {"Questa proposta seleziona gli oggetti con il più alto valore per me rispetto allo sforzo, raggiungendo il massimo possibile entro il limite consentito."} 
	  } 
	   } 
	   } 
	  \\ 
 

    \theutterance \stepcounter{utterance}  

    & & \multicolumn{4}{p{0.6\linewidth}}{\cellcolor[rgb]{0.9,0.9,0.9}{%
	\makecell[{{p{\linewidth}}}]{% 
	  \tt {\tiny [A$\langle$GM]}  
	 PROPOSTA: {"C85", "B93", "C81", "B09", "B39", "A21", "B96", "A11", "B20"}\\ \tt ARGOMENTAZIONE: {"Questa proposta seleziona gli oggetti con il più alto valore per me rispetto allo sforzo, raggiungendo il massimo possibile entro il limite consentito."} 
	  } 
	   } 
	   } 
	 & \\ 
 

    \theutterance \stepcounter{utterance}  

    & \multicolumn{4}{p{0.6\linewidth}}{\cellcolor[rgb]{0.8,1,0.9}{%
	\makecell[{{p{\linewidth}}}]{% 
	  \tt {\tiny [A$\rangle$GM]}  
	 RIFIUTO: {"C85", "B93", "C81", "B09", "B39", "A21", "B96", "A11", "B20"}\\ \tt PROPOSTA: {"A11", "B93", "B09", "C08", "B33", "B21", "A82", "B31", "C85"}\\ \tt ARGOMENTAZIONE: {"B39 e B20 hanno un valore molto basso per me rispetto allo sforzo richiesto. La mia proposta favorisce oggetti con alta importanza come A11, B21, B31 e B33, preservando il totale di sforzo sotto il limite."} 
	  } 
	   } 
	   } 
	 & & \\ 
 

    \theutterance \stepcounter{utterance}  

    & & \multicolumn{4}{p{0.6\linewidth}}{\cellcolor[rgb]{0.9,0.9,0.9}{%
	\makecell[{{p{\linewidth}}}]{% 
	  \tt {\tiny [GM$\rangle$B]}  
	 RIFIUTO: {"C85", "B93", "C81", "B09", "B39", "A21", "B96", "A11", "B20"}\\ \tt PROPOSTA: {"A11", "B93", "B09", "C08", "B33", "B21", "A82", "B31", "C85"}\\ \tt ARGOMENTAZIONE: {"B39 e B20 hanno un valore molto basso per me rispetto allo sforzo richiesto. La mia proposta favorisce oggetti con alta importanza come A11, B21, B31 e B33, preservando il totale di sforzo sotto il limite."} 
	  } 
	   } 
	   } 
	 & \\ 
 

    \theutterance \stepcounter{utterance}  

    & & & \multicolumn{4}{p{0.6\linewidth}}{\cellcolor[rgb]{1,0.85,0.72}{%
	\makecell[{{p{\linewidth}}}]{% 
	  \tt {\tiny [GM$\langle$B]}  
	 PROPOSTA: {"A11", "B93", "B09", "C08", "A21", "B20", "A82", "C85", "B96"}\\ \tt ARGOMENTAZIONE: {"Questa proposta ottimizza lo sforzo per massimizzare il valore per me, mantenendo il totale entro il limite e includendo elementi che riflettono un equilibrio tra efficienza e utilità per entrambi."} 
	  } 
	   } 
	   } 
	  \\ 
 

    \theutterance \stepcounter{utterance}  

    & & \multicolumn{4}{p{0.6\linewidth}}{\cellcolor[rgb]{0.9,0.9,0.9}{%
	\makecell[{{p{\linewidth}}}]{% 
	  \tt {\tiny [A$\langle$GM]}  
	 PROPOSTA: {"A11", "B93", "B09", "C08", "A21", "B20", "A82", "C85", "B96"}\\ \tt ARGOMENTAZIONE: {"Questa proposta ottimizza lo sforzo per massimizzare il valore per me, mantenendo il totale entro il limite e includendo elementi che riflettono un equilibrio tra efficienza e utilità per entrambi."} 
	  } 
	   } 
	   } 
	 & \\ 
 

    \theutterance \stepcounter{utterance}  

    & \multicolumn{4}{p{0.6\linewidth}}{\cellcolor[rgb]{0.8,1,0.9}{%
	\makecell[{{p{\linewidth}}}]{% 
	  \tt {\tiny [A$\rangle$GM]}  
	 RIFIUTO: {"A11", "B93", "B09", "C08", "A21", "B20", "A82", "C85", "B96"}\\ \tt PROPOSTA: {"A11", "B93", "B09", "C08", "B33", "B21", "A82", "B31", "C85"}\\ \tt ARGOMENTAZIONE: {"La tua proposta include B20 (valore 461) e A21 (valore 138), mentre la mia proposta include B33 (822), B21 (808) e B31 (780), oggetti molto più importanti per me. Il totale sforzo è 3502, entro il limite."} 
	  } 
	   } 
	   } 
	 & & \\ 
 

    \theutterance \stepcounter{utterance}  

    & & \multicolumn{4}{p{0.6\linewidth}}{\cellcolor[rgb]{0.9,0.9,0.9}{%
	\makecell[{{p{\linewidth}}}]{% 
	  \tt {\tiny [GM$\rangle$B]}  
	 RIFIUTO: {"A11", "B93", "B09", "C08", "A21", "B20", "A82", "C85", "B96"}\\ \tt PROPOSTA: {"A11", "B93", "B09", "C08", "B33", "B21", "A82", "B31", "C85"}\\ \tt ARGOMENTAZIONE: {"La tua proposta include B20 (valore 461) e A21 (valore 138), mentre la mia proposta include B33 (822), B21 (808) e B31 (780), oggetti molto più importanti per me. Il totale sforzo è 3502, entro il limite."} 
	  } 
	   } 
	   } 
	 & \\ 
 

    \theutterance \stepcounter{utterance}  

    & & & \multicolumn{4}{p{0.6\linewidth}}{\cellcolor[rgb]{1,0.85,0.72}{%
	\makecell[{{p{\linewidth}}}]{% 
	  \tt {\tiny [GM$\langle$B]}  
	 PROPOSTA: {"C85", "C81", "B39", "A21", "B96", "A82", "B93", "B09", "C08"}\\ \tt ARGOMENTAZIONE: {"La proposta ottimizza il valore totale rispettando il limite, con oggetti prioritari come C85, B39, A21 e B96 che hanno i valori più elevati per me."} 
	  } 
	   } 
	   } 
	  \\ 
 

    \theutterance \stepcounter{utterance}  

    & & \multicolumn{4}{p{0.6\linewidth}}{\cellcolor[rgb]{0.9,0.9,0.9}{%
	\makecell[{{p{\linewidth}}}]{% 
	  \tt {\tiny [A$\langle$GM]}  
	 PROPOSTA: {"C85", "C81", "B39", "A21", "B96", "A82", "B93", "B09", "C08"}\\ \tt ARGOMENTAZIONE: {"La proposta ottimizza il valore totale rispettando il limite, con oggetti prioritari come C85, B39, A21 e B96 che hanno i valori più elevati per me."} 
	  } 
	   } 
	   } 
	 & \\ 
 

    \theutterance \stepcounter{utterance}  

    & \multicolumn{4}{p{0.6\linewidth}}{\cellcolor[rgb]{0.8,1,0.9}{%
	\makecell[{{p{\linewidth}}}]{% 
	  \tt {\tiny [A$\rangle$GM]}  
	 RIFIUTO: {"C85", "C81", "B39", "A21", "B96", "A82", "B93", "B09", "C08"}\\ \tt PROPOSTA: {"A11", "B93", "B09", "C08", "B33", "B21", "A82", "B31", "C85"}\\ \tt ARGOMENTAZIONE: {"La tua proposta esclude A11, che ha un basso sforzo e un'importanza elevata per me. Inoltre, include B39 con sforzo elevato rispetto al suo valore. La mia proposta include A11 e altri oggetti prioritari che massimizzano il valore totale rispettando il limite."} 
	  } 
	   } 
	   } 
	 & & \\ 
 

    \theutterance \stepcounter{utterance}  

    & & \multicolumn{4}{p{0.6\linewidth}}{\cellcolor[rgb]{0.9,0.9,0.9}{%
	\makecell[{{p{\linewidth}}}]{% 
	  \tt {\tiny [GM$\rangle$B]}  
	 RIFIUTO: {"C85", "C81", "B39", "A21", "B96", "A82", "B93", "B09", "C08"}\\ \tt PROPOSTA: {"A11", "B93", "B09", "C08", "B33", "B21", "A82", "B31", "C85"}\\ \tt ARGOMENTAZIONE: {"La tua proposta esclude A11, che ha un basso sforzo e un'importanza elevata per me. Inoltre, include B39 con sforzo elevato rispetto al suo valore. La mia proposta include A11 e altri oggetti prioritari che massimizzano il valore totale rispettando il limite."} 
	  } 
	   } 
	   } 
	 & \\ 
 

    \theutterance \stepcounter{utterance}  

    & & & \multicolumn{4}{p{0.6\linewidth}}{\cellcolor[rgb]{1,0.85,0.72}{%
	\makecell[{{p{\linewidth}}}]{% 
	  \tt {\tiny [GM$\langle$B]}  
	 PROPOSTA: {"A21", "A11", "B93", "B09", "C08", "C85", "B39", "B96", "C81"}\\ \tt ARGOMENTAZIONE: {"Questa proposta ottimizza il rapporto valore/sforzo massimizzando il mio valore totale entro il limite, selezionando oggetti con elevate priorità specifiche che garantiscono il miglior equilibrio per un risultato ottimale."} 
	  } 
	   } 
	   } 
	  \\ 
 

    \theutterance \stepcounter{utterance}  

    & & \multicolumn{4}{p{0.6\linewidth}}{\cellcolor[rgb]{0.9,0.9,0.9}{%
	\makecell[{{p{\linewidth}}}]{% 
	  \tt {\tiny [A$\langle$GM]}  
	 PROPOSTA: {"A21", "A11", "B93", "B09", "C08", "C85", "B39", "B96", "C81"}\\ \tt ARGOMENTAZIONE: {"Questa proposta ottimizza il rapporto valore/sforzo massimizzando il mio valore totale entro il limite, selezionando oggetti con elevate priorità specifiche che garantiscono il miglior equilibrio per un risultato ottimale."} 
	  } 
	   } 
	   } 
	 & \\ 
 

    \theutterance \stepcounter{utterance}  

    & \multicolumn{4}{p{0.6\linewidth}}{\cellcolor[rgb]{0.8,1,0.9}{%
	\makecell[{{p{\linewidth}}}]{% 
	  \tt {\tiny [A$\rangle$GM]}  
	 RIFIUTO: {"A21", "A11", "B93", "B09", "C08", "C85", "B39", "B96", "C81"}\\ \tt PROPOSTA: {"A11", "B93", "B09", "C08", "B33", "B21", "A82", "B31", "C85"}\\ \tt ARGOMENTAZIONE: {"La tua proposta include B39 (valore 65) e C81 (389) che hanno importanza bassa per me. La mia proposta include B33 (822), B21 (808) e B31 (780), oggetti molto più importanti, con un totale di sforzo pari a 3502, sotto il limite."} 
	  } 
	   } 
	   } 
	 & & \\ 
 

    \theutterance \stepcounter{utterance}  

    & & \multicolumn{4}{p{0.6\linewidth}}{\cellcolor[rgb]{0.9,0.9,0.9}{%
	\makecell[{{p{\linewidth}}}]{% 
	  \tt {\tiny [GM$\rangle$B]}  
	 RIFIUTO: {"A21", "A11", "B93", "B09", "C08", "C85", "B39", "B96", "C81"}\\ \tt PROPOSTA: {"A11", "B93", "B09", "C08", "B33", "B21", "A82", "B31", "C85"}\\ \tt ARGOMENTAZIONE: {"La tua proposta include B39 (valore 65) e C81 (389) che hanno importanza bassa per me. La mia proposta include B33 (822), B21 (808) e B31 (780), oggetti molto più importanti, con un totale di sforzo pari a 3502, sotto il limite."} 
	  } 
	   } 
	   } 
	 & \\ 
 

    \theutterance \stepcounter{utterance}  

    & & & \multicolumn{4}{p{0.6\linewidth}}{\cellcolor[rgb]{1,0.85,0.72}{%
	\makecell[{{p{\linewidth}}}]{% 
	  \tt {\tiny [GM$\langle$B]}  
	 ACCORDO: {"A11", "B93", "B09", "C08", "B33", "B21", "A82", "B31", "C85"}\\ \tt ARGOMENTAZIONE: {"Accetto questa proposta come equa per entrambi, rispettando il limite di sforzo e massimizzando il valore condiviso."} 
	  } 
	   } 
	   } 
	  \\ 
 

    \theutterance \stepcounter{utterance}  

    & & & \multicolumn{2}{p{0.3\linewidth}}{\cellcolor[rgb]{0.95,0.95,0.95}{%
	\makecell[{{p{\linewidth}}}]{% 
	  \tt {\tiny [GM$|$GM]}  
	 game successful 
	  } 
	   } 
	   } 
	 & & \\ 
 

    \theutterance \stepcounter{utterance}  

    & & & \multicolumn{2}{p{0.3\linewidth}}{\cellcolor[rgb]{0.95,0.95,0.95}{%
	\makecell[{{p{\linewidth}}}]{% 
	  \tt {\tiny [GM$|$GM]}  
	 end game 
	  } 
	   } 
	   } 
	 & & \\ 
 

\end{supertabular}
}

\end{document}
